\documentclass[leqno]{article}
% Shared mark scheme preamble for resources.danieljsmith.org: the pack
% preamble, unchanged, then the scheme's own layer. A scheme's .tex is
% rendered from its YAML (scripts/README.md, Mark schemes), never edited.
\input{\CurrentFilePath/pack_preamble.tex}

% The page: equal side margins, so the table runs the full width of the
% header rule. The footers' right edge follows the text's.
\geometry{left=0.8in, right=0.8in}
\makeatletter
\renewcommand{\djs@footR}[1]{%
  \put(\dimexpr1in+\hoffset+\oddsidemargin+\textwidth\relax,
    -\dimexpr\paperheight-0.32in\relax){\makebox[0pt][r]{#1}}}
\makeatother

\usepackage[export]{adjustbox}
\usepackage{tabularray}
\definecolor{msGrid}{RGB}{176, 172, 164}
\definecolor{msTint}{RGB}{245, 243, 239}

% The header on every page, the front page included, and the booklet label
% under the contents-page title. \djsQuestionPageStyle stays empty.
\newcommand{\djsSchemeHeader}[1]{%
  \djsPackHeader[markscheme]{\textit{Mark scheme} - #1}%
  \renewcommand{\djsPackTitle}{#1}%
  \renewcommand{\djsBookletLabel}{Mark scheme}%
}

% The front page: the contents heading, spaced as tocloft spaces the packs'
% "Questions", then the contents table. A question's name and page number
% both link to its first page.
\newcommand{\msContentsHeading}{%
  \vskip\cftbeforetoctitleskip
  {\noindent\cfttoctitlefont\contentsname\par}%
  \vskip\cftaftertoctitleskip}
\NewTblrEnviron{mscontents}
\SetTblrInner[mscontents]{
  colspec = {Q[wd=26mm, l] X[l] Q[wd=16mm, c] Q[wd=14mm, r]},
  width = \linewidth,
  rowsep = 4pt, colsep = 6pt,
  hlines = {0.4pt, msGrid}, vlines = {0.4pt, msGrid},
  hline{2} = {2.2pt, scarlet},
  row{1} = {font=\small\bfseries},
  row{Z} = {font=\bfseries},
  column{2} = {font=\small},
}
\newcommand{\msContentsQuestion}[1]{\hyperref[q:#1]{Question #1}}
\newcommand{\msContentsPage}[1]{\hyperref[q:#1]{\pageref*{q:#1}}}
\newcommand{\msSep}{\hspace{0.9em}}

% The conventions pages' layout: headings in the documents' own bold, tables
% ruled and padded as the scheme's, roomy lists.
\newcommand{\msHead}[1]{%
  \par\vspace{18pt}%
  {\noindent\bfseries #1\par}%
  \vspace{6pt}}
\NewTblrEnviron{mstblr}
\SetTblrInner[mstblr]{rowsep = 6pt, colsep = 6pt,
  hlines = {0.4pt, msGrid}, vlines = {0.4pt, msGrid}}
\newenvironment{msList}
  {\begin{itemize}[leftmargin=1.2em, itemsep=5pt, topsep=3pt, parsep=0pt]}
  {\end{itemize}}

% A question's first page: the link target and the PDF bookmark, as
% \questionitem gives the packs.
\newcommand{\msQnum}{}
\newcommand{\msQuestion}[1]{%
  \clearpage
  \renewcommand{\msQnum}{#1}%
  \phantomsection\label{q:#1}%
  \addcontentsline{toc}{subsection}{\protect\djsTocEntry{#1}}}

% The scheme's table: Question 16mm, Scheme the rest, Marks 16mm, Notes
% 44mm; every cell ruled in msGrid; the scarlet rule under the head row. A
% long table has no caption or "continued" foot; a continuation page opens
% with a muted line above the repeated head row. The rules across a part's
% rows, the alternative's tint and the total row are the rendered table's.
\NewTblrTheme{ms}{%
  \DefTblrTemplate{firsthead}{default}{}%
  \DefTblrTemplate{middlehead,lasthead}{default}{%
    {\small\itshape\color{djsMuted}Question \msQnum\ continued\par}%
    \vspace{6pt}}%
  \DefTblrTemplate{firstfoot,middlefoot,lastfoot}{default}{}%
}
\NewTblrEnviron{msscheme}
\SetTblrOuter[msscheme]{long, theme=ms, entry=none, label=none}
\SetTblrInner[msscheme]{
  colspec = {Q[wd=16mm, l] X[l] Q[wd=16mm, c] Q[wd=44mm, l]},
  width = \linewidth,
  rowsep = 6pt, colsep = 6pt,
  rowhead = 1,
  vlines = {0.4pt, msGrid},
  hline{1} = {0.4pt, msGrid}, hline{2} = {2.2pt, scarlet},
  hline{Z} = {0.4pt, msGrid},
  column{1} = {font=\bfseries}, column{3} = {font=\bfseries},
  column{4} = {font=\small},
  row{1} = {rowsep=4pt, font=\small\bfseries},
}

% Two marks on a row: each note starts with its code.
\newcommand{\msNote}[2]{\textbf{#1}\enspace #2}
% A full-row figure: the worked solution's own picture, scaled uniformly and
% only down, to the table less the Question column. The picture is drawn
% with the line width it has in the solutions document, at its list depth
% there (the optional argument, 1 by default: a part's working), so one
% sized by \linewidth keeps its size; tabularray sets the cell's \linewidth
% to the cell's measured width, which would shrink it at every pass. The
% widths are measured once, in a solution box at the packs' text width
% (pack_preamble.tex's margins, 0.8in and 1in).
\newlength{\msSpanW}
\makeatletter
\newcommand{\ms@keepLW}[1]{%
  \expandafter\xdef\csname ms@solnLW@#1\endcsname{\the\linewidth}}
\AtBeginDocument{%
  \setlength{\msSpanW}{\dimexpr\textwidth-16mm-28pt\relax}%
  \setbox\z@\vbox{%
    \begin{minipage}{\dimexpr\paperwidth-1.8in\relax}
      \begin{djsSolution}
        \ms@keepLW{0}%
        \begin{enumerate}\item \ms@keepLW{1}%
        \begin{enumerate}\item \ms@keepLW{2}%
        \begin{enumerate}\item \ms@keepLW{3}%
        \end{enumerate}\end{enumerate}\end{enumerate}
      \end{djsSolution}
    \end{minipage}}}
\newcommand{\msFigure}[2][1]{%
  \adjustbox{max width=\msSpanW, max totalheight=0.6\textheight, valign=t}{%
    \hsize=\csname ms@solnLW@#1\endcsname
    \linewidth=\hsize
    #2}}
\makeatother
% A figure mark's feature, in the Scheme column.
\newcommand{\msFeature}[1]{\textit{#1}}
% An alternative's title row.
\newcommand{\msAlternative}[1]{\textit{Alternative:} #1}
% Drafts only: a queried mark's code, the query's why under its note, and a
% flag in its part's place.
\newcommand{\msQueried}[2]{#1~\(\to\) #2?}
\newcommand{\msQuery}[1]{{\itshape\color{djsMuted}Query: #1}}
\newcommand{\msFlag}[2]{\textbf{#1:}\enspace #2}


\djsSchemeHeader{Dimensional Analysis}
\djsSitePath{Further Mechanics}{Miscellaneous Mechanics}

\djsLastUpdated{30 October 2026}
\begin{document}
\djsSetQuestionCount{10}
\djsAddMarks{1}{5}
\djsAddMarks{2}{5}
\djsAddMarks{3}{6}
\djsAddMarks{4}{4}
\djsAddMarks{5}{6}
\djsAddMarks{6}{6}
\djsAddMarks{7}{2}
\djsAddMarks{8}{4}
\djsAddMarks{9}{6}
\djsAddMarks{10}{4}
\vspace*{6pt}
\djsContentsTitle
\vspace{30pt}
\msContentsHeading
\begin{mscontents}{}
Question & Marks by part & Marks & Page \\
\msContentsQuestion{1} & (a)\,1\msSep (b)\,4 & 5 & \msContentsPage{1} \\
\msContentsQuestion{2} & (a)\,1\msSep (b)\,4 & 5 & \msContentsPage{2} \\
\msContentsQuestion{3} & (a)\,2\msSep (b)\,2\msSep (c)\,2 & 6 & \msContentsPage{3} \\
\msContentsQuestion{4} & (a)\,1\msSep (b)\,3 & 4 & \msContentsPage{4} \\
\msContentsQuestion{5} & (a)\,1\msSep (b)\,1\msSep (c)\,4 & 6 & \msContentsPage{5} \\
\msContentsQuestion{6} & & 6 & \msContentsPage{6} \\
\msContentsQuestion{7} & & 2 & \msContentsPage{7} \\
\msContentsQuestion{8} & (a)\,2\msSep (b)\,2 & 4 & \msContentsPage{8} \\
\msContentsQuestion{9} & (a)\,2\msSep (b)\,4 & 6 & \msContentsPage{9} \\
\msContentsQuestion{10} & & 4 & \msContentsPage{10} \\
\SetCell[c=2]{r} Total & & 48 & \\
\end{mscontents}
\newpage
% The conventions pages, shared by every pack's mark scheme: one source of
% truth, \input after the front page, with no strand or pack logic. The keys
% on the first side, the rules on the second. The scheme layer defines
% \msHead, msList and mstblr.
\vspace*{6pt}%
{\noindent\cfttoctitlefont Reading the mark scheme\par}%
\vspace{10pt}%
These rules apply throughout, unless a note says otherwise.

\msHead{Marks}
Every mark is worth one, so \textbf{M1 A1} on a row is two marks. Only marks
coded ft follow through.\par
\vspace{10pt}%
\begin{mstblr}{
  colspec = {Q[wd=15mm, l] X[l]},
  width = \linewidth,
  hline{2} = {2.2pt, scarlet},
  column{1} = {font=\bfseries},
  row{1} = {font=\small\bfseries, rowsep=4pt},
}
Code & Meaning \\
M1 & Method: a correct method, applied to this question. A slip doesn't cost it if the method is clear. A formula must be used, not just quoted. \\
dM1 & Dependent method: as M1, but only if the method mark before it was earned, unless its note names another. \\
A1 & Accuracy: a correct answer or step. It needs the method mark it follows. \\
A1ft & As A1, following through from an earlier answer that the note names. \\
A1~AG & The last mark of a part whose answer is printed in the question: complete working, with no errors. \\
B1 & A correct value, fact or statement. It needs no method mark. \\
B1ft & As B1, following through from an earlier answer that the note names. \\
E1 & A written explanation, interpretation or conclusion. Concluding a proof by induction needs every earlier mark. \\
\end{mstblr}

\msHead{Abbreviations}
ft and AG qualify a code in the Marks column. The others appear in the
notes.\par
\vspace{10pt}%
\begin{mstblr}{
  colspec = {Q[wd=17mm, l] X[l] Q[wd=15mm, l] X[l]},
  width = \linewidth,
  column{1} = {font=\bfseries}, column{3} = {font=\bfseries},
}
ft & follow through & AG & answer given in the question \\
cao & correct answer only, with no follow-through & cso & correct solution only: no errors in the part \\
oe & or an equivalent form & awrt & any answer that rounds to this value \\
soi & seen, or implied by later working & isw & ignore working after a correct answer \\
SC & special case: credit outside the main method & their & the student's own earlier value \\
condone & a slip that doesn't cost the mark & [\;\,] & working the mark doesn't need \\
s.f. & significant figures & d.p. & decimal places \\
\end{mstblr}

\newpage
\msHead{Reading the scheme}
\begin{msList}
\item Each row is a line of working, with the marks it earns beside it.
  Where a row carries two marks, each note starts with its code.
\item \textbf{[\textit{n}]} ends each part with its total.
\item A shaded \textit{Alternative} is a second method for the same total.
  Mark an answer by the method it follows.
\item Any other valid method earns the equivalent marks, unless the question
  names a method or says \textit{hence}.
\item A sketch earns one B1 for each feature named below its figure.
\end{msList}

\msHead{Accuracy and answers}
\begin{msList}
\item Unless a question says otherwise, give answers to 3 significant
  figures, and angles in degrees to 1 decimal place. Angles in radians take
  3 significant figures.
\item A final answer with more figures is accepted if it rounds to the
  scheme's value; one with fewer is not. An answer spoiled by rounding too
  early loses its accuracy mark.
\item Where a question asks for an exact answer, a decimal loses the final
  accuracy mark. Otherwise either is accepted.
\item Wrong or missing units cost nothing unless the question asks for them.
\item A correct answer with no working earns full marks, unless the question
  needs the working: show that, prove, hence, a method it names, or an exact
  answer. A wrong answer with no working earns nothing. Calculator answers
  are treated the same way.
\item A correct answer, once reached, isn't spoiled by later working (isw).
\end{msList}

\msHead{Value of $\boldsymbol{g}$}
\begin{msList}
\item Take $g = 9.8\,\mathrm{m\,s^{-2}}$ unless a question states a value,
  which is then the only one accepted. Answers that use $g = 9.8$ may be
  given to 2 or 3 significant figures.
\item Where a question doesn't state $g$, an answer worked with $g = 10$
  throughout earns the same marks, keeping any value printed in the
  question. The worked solutions list the final answers with $g = 10$.
\end{msList}


\msQuestion{1}
\begin{msscheme}{
  hline{3,5-8} = {2-Z}{0.4pt, msGrid},
  hline{4,9} = {0.4pt, msGrid},
  row{9} = {font=\bfseries},
}
Question & Scheme & Marks & Notes \\
1(a) & {\([F] = \mathrm{MLT^{-2}}\)} & B1 & {Not for the units \(\mathrm{kg\,m\,s^{-2}}\) or \(\mathrm{N}\).} \\
\nopagebreak & & [1] & \\
1(b) & {\(\mathrm{L} = \mathrm{M}^{\alpha}(\mathrm{LT^{-1}})^{\beta}(\mathrm{MLT^{-2}})^{\gamma}\)} & M1 & {Equates the dimensions of both sides, using their \([F]\) from (a) for \([R]\). Condone one slip in a dimension. Powers need not be expanded.} \\
& {M: \(\alpha + \gamma = 0\)\\[4pt] L: \(\beta + \gamma = 1\)\\[4pt] T: \(-\beta - 2\gamma = 0\)} & A1 & {All three correct, oe.} \\
& {\(\Rightarrow \alpha = \ldots,\ \beta = \ldots,\ \gamma = \ldots\)} & dM1 & {Solves their three equations by any means to find all three powers.} \\
& {\(\alpha = 1,\ \beta = 2,\ \gamma = -1\)} & A1 & \\
\nopagebreak & & [4] & \\
\nopagebreak \SetCell[c=2]{r} Total for Question 1 & & 5 & \\
\end{msscheme}

\msQuestion{2}
\begin{msscheme}{
  hline{3,5-8} = {2-Z}{0.4pt, msGrid},
  hline{4,9} = {0.4pt, msGrid},
  row{9} = {font=\bfseries},
}
Question & Scheme & Marks & Notes \\
2(a) & {\([G] = \mathrm{M^{-1}L^3T^{-2}}\)} & B1 & {Symbols in any order. Not for the units \(\mathrm{m^3\,kg^{-1}\,s^{-2}}\).} \\
\nopagebreak & & [1] & \\
2(b) & {\(\mathrm{T} = (\mathrm{M^{-1}L^3T^{-2}})^{\alpha}\mathrm{M}^{\beta}\mathrm{L}^{\gamma}\)} & M1 & {Forms a dimensional equation using their \([G]\) from (a) and the dimensions of \(T\), \(M\) and \(r\), each raised to its power. Condone one slip in a dimension.} \\
& {M: \(-\alpha + \beta = 0\)\\[4pt] L: \(3\alpha + \gamma = 0\)\\[4pt] T: \(-2\alpha = 1\)} & A1 & {All three correct, oe.} \\
& {\(\alpha = \ldots,\ \beta = \ldots,\ \gamma = \ldots\)} & dM1 & {Solves their three equations from equating powers for all three values. Condone sign slips.} \\
& {\(\alpha = -\tfrac12,\ \beta = -\tfrac12,\ \gamma = \tfrac32\)} & A1 & {All three correct, oe.} \\
\nopagebreak & & [4] & \\
\nopagebreak \SetCell[c=2]{r} Total for Question 2 & & 5 & \\
\end{msscheme}

\msQuestion{3}
\begin{msscheme}{
  hline{3-4,6-7,9-10} = {2-Z}{0.4pt, msGrid},
  hline{5,8,11} = {0.4pt, msGrid},
  row{11} = {font=\bfseries},
}
Question & Scheme & Marks & Notes \\
3(a) & {\([\Gamma] = [g][r^2] = \mathrm{LT^{-2}} \times \mathrm{L^2}\)} & M1 & {Uses \(\Gamma = gr^2\) with \([g] = \mathrm{LT^{-2}}\) and \([r] = \mathrm{L}\). Condone a slip in combining powers.} \\
& {\([\Gamma] = \mathrm{L^3T^{-2}}\)} & A1 & \\
\nopagebreak & & [2] & \\
3(b) & {\([v^2] = [u^2] = \mathrm{L^2T^{-2}}\)\\[4pt] \(\left[2\Gamma\left(\frac{1}{R}-\frac{1}{r}\right)\right] = \their{\mathrm{L^3T^{-2}}} \times \mathrm{L^{-1}}\)} & M1 & {Finds the dimensions of the speed terms and of the \(\Gamma\) term, using their \([\Gamma]\) of the form \(\mathrm{L}^\alpha\mathrm{T}^\beta\) and \(\mathrm{L^{-1}}\) for the bracket.} \\
& {\(\phantom{\left[2\Gamma\right]} = \mathrm{L^2T^{-2}}\)\\[4pt] All terms \(\mathrm{L^2T^{-2}}\), so consistent} & A1~AG & {cso, from a correct \([\Gamma]\). \([\Gamma]\) must not come from assuming this equation consistent. All three terms shown as \(\mathrm{L^2T^{-2}}\), with a conclusion in words. Comment on the factor 2 not required.} \\
\nopagebreak & & [2] & \\
3(c) & {\([W] = \mathrm{M} \times \their{\mathrm{L^3T^{-2}}} \times \mathrm{L^{-1}}\)} & M1 & {Multiplies \([m] = \mathrm{M}\), their \([\Gamma]\) of the form \(\mathrm{L}^\alpha\mathrm{T}^\beta\) and \(\mathrm{L^{-1}}\) for the bracket. Condone a slip in combining powers.} \\
& {\([W] = \mathrm{ML^2T^{-2}}\); \(\mathrm{kg\,m^2\,s^{-2}}\) or J} & A1 & {Both the dimensions and the units; accept J, \(\mathrm{N\,m}\) or \(\mathrm{kg\,m^2\,s^{-2}}\).} \\
\nopagebreak & & [2] & \\
\nopagebreak \SetCell[c=2]{r} Total for Question 3 & & 6 & \\
\end{msscheme}

\msQuestion{4}
\begin{msscheme}{
  hline{3,5-7} = {2-Z}{0.4pt, msGrid},
  hline{4,8} = {0.4pt, msGrid},
  row{8} = {font=\bfseries},
}
Question & Scheme & Marks & Notes \\
4(a) & {\([p] = \mathrm{ML^{-1}T^{-2}}\)} & B1 & {Not for the units \(\mathrm{Pa}\) or \(\mathrm{kg\,m^{-1}\,s^{-2}}\).} \\
\nopagebreak & & [1] & \\
4(b) & {\(\mathrm{ML^{-1}T^{-2}} = (\mathrm{ML^{-3}})^{\alpha}(\mathrm{L})^{\beta}(\mathrm{LT^{-2}})\)\\[4pt] \(\phantom{\mathrm{ML^{-1}T^{-2}}} = \mathrm{M}^{\alpha}\mathrm{L}^{-3\alpha+\beta+1}\mathrm{T^{-2}}\)} & M1 & {Equates their \([p]\) to a product of powers of dimensions for \(\rho\), \(h\) and \(g\), with \(k\) dimensionless. Condone one slip in a dimension.} \\
& {M: \(\alpha = 1\)} & A1 & \\
& {L: \(-3\alpha + \beta + 1 = -1 \Rightarrow \beta = 1\)} & A1 & \\
\nopagebreak & & [3] & \\
\nopagebreak \SetCell[c=2]{r} Total for Question 4 & & 4 & \\
\end{msscheme}

\msQuestion{5}
\begin{msscheme}{
  hline{3,5,7-10} = {2-Z}{0.4pt, msGrid},
  hline{4,6,11} = {0.4pt, msGrid},
  row{11} = {font=\bfseries},
}
Question & Scheme & Marks & Notes \\
5(a) & {\([\rho] = \mathrm{ML^{-3}}\)} & B1 & {Accept \(\mathrm{M/L^3}\) oe. Not for the units \(\mathrm{kg\,m^{-3}}\).} \\
\nopagebreak & & [1] & \\
5(b) & {Not possible: \(\tfrac12\) is dimensionless} & E1 & {States it is not possible because the constant is dimensionless, or a pure number. Not for ``no'' alone.} \\
\nopagebreak & & [1] & \\
5(c) & {\(\mathrm{MLT^{-2}} = (\mathrm{ML^{-3}})^{\alpha}(\mathrm{L^2})^{\beta}(\mathrm{LT^{-1}})^{\gamma}\)} & M1 & {Forms a dimensional equation using their \([\rho]\) from (a) and the dimensions of \(F\), \(A\) and \(v\). Condone one slip in a dimension, and \(\frac12\) left in.} \\
& {M: \(\alpha = 1\)\\[4pt] T: \(-\gamma = -2\)\\[4pt] L: \(-3\alpha + 2\beta + \gamma = 1\)} & A1 & {All three equations correct, oe.} \\
& {\(-3(1) + 2\beta + 2 = 1 \Rightarrow \beta = \ldots\)} & dM1 & {Solves their three equations, substituting their \(\alpha\) and \(\gamma\) to find \(\beta\).} \\
& {\(\alpha = 1\), \(\beta = 1\), \(\gamma = 2\)} & A1 & \\
\nopagebreak & & [4] & \\
\nopagebreak \SetCell[c=2]{r} Total for Question 5 & & 6 & \\
\end{msscheme}

\msQuestion{6}
\begin{msscheme}{
  hline{3-7} = {2-Z}{0.4pt, msGrid},
  hline{8} = {0.4pt, msGrid},
  row{8} = {font=\bfseries},
}
Question & Scheme & Marks & Notes \\
6 & {\(\left[k_1 mv\frac{\mathrm{d}v}{\mathrm{d}x}\right] = \mathrm{M}\times\mathrm{LT^{-1}}\times\frac{\mathrm{LT^{-1}}}{\mathrm{L}}\)} & M1 & {Replaces the quantities in at least one term by their dimensions and combines them. Condone one slip.} \\
& {\(\phantom{\left[k_1 mv\frac{\mathrm{d}v}{\mathrm{d}x}\right]} = \mathrm{MLT^{-2}}\)} & A1 & {Accept \(v\frac{\mathrm{d}v}{\mathrm{d}x}\) treated as an acceleration, \(\mathrm{LT^{-2}}\).} \\
& {\([k_2\mu mg] = \mathrm{M}\times\mathrm{LT^{-2}} = \mathrm{MLT^{-2}}\)} & A1 & {\(k_2\) and \(\mu\) dimensionless, stated or implied by their omission.} \\
& {\([Q] = \mathrm{ML^2T^{-2}}\)} & B1 & \\
& {\(\left[k_3\frac{Q}{x}\right] = \frac{\mathrm{ML^2T^{-2}}}{\mathrm{L}} = \mathrm{MLT^{-2}}\)} & A1 & \\
& {\([F] = \mathrm{MLT^{-2}}\), the same as every term\\[4pt] \(\Rightarrow\) dimensionally consistent} & E1 & {States \([F] = \mathrm{MLT^{-2}}\) and concludes consistent, with all three terms correctly shown as \(\mathrm{MLT^{-2}}\).} \\
\nopagebreak \SetCell[c=2]{r} Total for Question 6 & & 6 & \\
\end{msscheme}

\msQuestion{7}
\begin{msscheme}{
  hline{3} = {2-Z}{0.4pt, msGrid},
  hline{4} = {0.4pt, msGrid},
  row{4} = {font=\bfseries},
}
Question & Scheme & Marks & Notes \\
7 & {Exponent dimensionless: \(\left[\frac{kt}{m}\right] = 1\)\\[4pt] \(\Rightarrow \frac{[k]\mathrm{T}}{\mathrm{M}} = 1 \Rightarrow [k] = \ldots\)} & M1 & {Sets the dimensions of \(\frac{kt}{m}\) equal to 1 with \([t]\) and \([m]\) substituted, or equates \(\left[\frac{mg}{k}\right]\) with \([v]\) in M, L, T.} \\
& {\([k] = \mathrm{MT^{-1}}\)} & A1 & {Any equivalent form. Not for the units \(\mathrm{kg\,s^{-1}}\).} \\
\nopagebreak \SetCell[c=2]{r} Total for Question 7 & & 2 & \\
\end{msscheme}

\msQuestion{8}
\begin{msscheme}{
  hline{3-4,6-7} = {2-Z}{0.4pt, msGrid},
  hline{5,8} = {0.4pt, msGrid},
  row{8} = {font=\bfseries},
}
Question & Scheme & Marks & Notes \\
8(a) & {\([T]^2 = \frac{[a]^3}{[\mu]} \Rightarrow \mathrm{T^2} = \frac{\mathrm{L^3}}{[\mu]}\)} & M1 & {Substitutes \([T]=\mathrm{T}\) and \([a]=\mathrm{L}\), with \(4\pi^2\) treated as dimensionless, into the equation. Condone slips in rearranging.} \\
& {\([\mu] = \mathrm{L^3T^{-2}}\)} & A1 & {oe, such as \(\mathrm{L^3/T^2}\). Not for the units \(\mathrm{m^3\,s^{-2}}\).} \\
\nopagebreak & & [2] & \\
8(b) & {\([H] = \mathrm{M}\sqrt{\their{\mathrm{L^3T^{-2}}}\,\mathrm{L}}\)\\[4pt] \(\phantom{[H]} = \mathrm{M}\left(\mathrm{L^4T^{-2}}\right)^{\frac12}\)} & M1 & {Uses \([m]=\mathrm{M}\), \([a]=\mathrm{L}\), \(1-e^2\) dimensionless and their \([\mu]\) in terms of \(\mathrm{L}\) and \(\mathrm{T}\), and halves the powers under the root.} \\
& {\([H] = \mathrm{ML^2T^{-1}}\)} & A1 & {oe. Not for the units \(\mathrm{kg\,m^2\,s^{-1}}\).} \\
\nopagebreak & & [2] & \\
\nopagebreak \SetCell[c=2]{r} Total for Question 8 & & 4 & \\
\end{msscheme}

\msQuestion{9}
\begin{msscheme}{
  hline{3-4,6-9} = {2-Z}{0.4pt, msGrid},
  hline{5,10} = {0.4pt, msGrid},
  row{10} = {font=\bfseries},
}
Question & Scheme & Marks & Notes \\
9(a) & {\([k][v]^2 = [F]\)\\[4pt] \(\Rightarrow [k](\mathrm{LT^{-1}})^2 = \mathrm{MLT^{-2}}\)} & M1 & {Equates the dimensions of \(kv^2\) to those of force, with \([F]\) in terms of \(\mathrm{M}\), \(\mathrm{L}\) and \(\mathrm{T}\). Condone a slip in squaring \([v]\).} \\
& {\([k] = \mathrm{ML^{-1}}\)} & A1 & {oe, such as \(\mathrm{ML^{-1}T^0}\). Not for the units \(\mathrm{kg\,m^{-1}}\).} \\
\nopagebreak & & [2] & \\
9(b) & {\([k^2 m^\alpha u^\beta g^\gamma] = \mathrm{L}\)\\[4pt] \((\their{\mathrm{ML^{-1}}})^2\,\mathrm{M}^\alpha(\mathrm{LT^{-1}})^\beta(\mathrm{LT^{-2}})^\gamma = \mathrm{L}\)} & M1 & {Sets the dimensions of the third term equal to \(\mathrm{L}\), or to those of another term, using their \([k]\) in \(\mathrm{M}\), \(\mathrm{L}\), \(\mathrm{T}\). Condone one slip in \([u]\) or \([g]\).} \\
& {M: \(2 + \alpha = 0\)\\[4pt] L: \(-2 + \beta + \gamma = 1\)\\[4pt] T: \(-\beta - 2\gamma = 0\)} & A1 & {All three equations correct, oe.} \\
& {\(\beta = -2\gamma \Rightarrow -2 - 2\gamma + \gamma = 1\)\\[4pt] \(\phantom{\beta = -2\gamma} \Rightarrow \gamma = \ldots\)} & dM1 & {Solves their three equations to find all of \(\alpha\), \(\beta\) and \(\gamma\), by any method.} \\
& {\(\alpha = -2,\ \beta = 6,\ \gamma = -3\)} & A1 & \\
\nopagebreak & & [4] & \\
\nopagebreak \SetCell[c=2]{r} Total for Question 9 & & 6 & \\
\end{msscheme}

\msQuestion{10}
\begin{msscheme}{
  hline{3-5} = {2-Z}{0.4pt, msGrid},
  hline{6} = {0.4pt, msGrid},
  row{6} = {font=\bfseries},
}
Question & Scheme & Marks & Notes \\
10 & {\((\mathrm{LT^{-1}})^\delta = \mathrm{L^2T^{-2}}\)\\[4pt] \((\mathrm{LT^{-1}})^\gamma = \sqrt{\mathrm{L^2T^{-2}}} = \mathrm{LT^{-1}}\)} & M1 & {Equates the dimensions of the terms added inside the square root, or of \(u^\gamma\sin\theta\) with the root, taking \(\sin\theta\) as dimensionless. Condone one dimension slip.} \\
& {\(\delta = 2\), \(\gamma = 1\)} & A1 & \\
& {\(\mathrm{L} = \frac{(\mathrm{LT^{-1}})^\alpha}{(\mathrm{LT^{-2}})^\beta}\,\mathrm{LT^{-1}}\)\\[4pt] \(\Rightarrow \alpha - \beta + 1 = 1,\ -\alpha + 2\beta - 1 = 0\)\\[4pt] \(\Rightarrow \alpha = \ldots,\ \beta = \ldots\)} & M1 & {Uses their bracket dimensions as a power product of \(\mathrm{L}\) and \(\mathrm{T}\), equates powers of \(\mathrm{L}\) and \(\mathrm{T}\) with \([R]\) and solves for \(\alpha\) and \(\beta\). Condone sign slips.} \\
& {\(\alpha = 1\), \(\beta = 1\)} & A1 & \\
\nopagebreak \SetCell[c=2]{r} Total for Question 10 & & 4 & \\
\end{msscheme}

\end{document}
